\selectlanguage{ngerman}%
\frenchspacing%
\begin{ledgroupsized}[r]{120mm}%
\footnotesize%
\pstart%
\noindent%
\textbf{Überlieferung:}%
\pend%
\end{ledgroupsized}%
\begin{ledgroupsized}[r]{114mm}%
\footnotesize%
\pstart%
\parindent -6mm%
\makebox[6mm][l]{\textit{E}}%
Druck nach unbekannter Vorlage:
\textsc{Gerland} 1906, S.\;5
(Nr.\;3 mit der Überschrift \textit{Wesen der Luft}).%
\cite{00197}
Die verschollene handschriftliche Vorlage war laut \textit{E} ein \textit{Blatt von der Hand Leibnizens, der rechts darauf "Cassini"\indexN{\textls{Cassini}, Giovanni Domenico 1625\textendash1712} gesetzt hat}.
Der oft fehlerhafte Text von \textit{E} (Gerlands Transkription, die wohl nur den Ausschnitt eines umfangreicheren Textes reproduziert) wird hier wiedergegeben und gegebenenfalls berichtigt oder kommentiert;
die Rechtschreibung wird an die Editionsregeln angepasst.
\pend%
\end{ledgroupsized}%
\vspace{5mm}%
\begin{ledgroup}%
\footnotesize%
\pstart%
\noindent%
\textbf{Datierungsgründe:}
Die undatierte Bemerkung N.\;\ref{56214}, wohl ein Ausschnitt aus einem umfangreicheren, nicht ermittelten Text Leibnizens, ist nach heutigem Wissensstand lediglich im Druck \textit{E} überliefert.
Befasst ist sie mit dem Gedanken, dass die atmosphärische Luft sich aus zwei funktional verscheinenen Bestandteilen zusammensetze: Von dem einen würde die Refraktion der Lichtstrahle, von dem anderen die Elastizität und Schwere der Luft selbst herrühren.
Hierbei knüpft Leibniz an Beobachtungen an, die Gian Domenico Cassini\indexN{\textls{Cassini}, Giovanni Domenico 1625\textendash1712} bereits im Jahr 1672 durchgeführt, aber wahrscheinlich erst in dem 1693 gedruckten Sammelband \textit{Recueil d’observations faites en plusieurs voyages par ordre de sa majesté, pour perfectionner l’astronomie et la géographie} veröffentlicht hatte (siehe unten den Erläuterungsapparat, \refpassage{Gerland_1906_005b_Cassini_wmc-1}{Gerland_1906_005b_Cassini_wmc-2}).\cite{02605}
Spätestens 1701 hatte Leibniz diese Veröffentlichung zur Kenntnis genommen, wie sein Brief an Joachim Tiede\indexN{\textls{Tiede}, Joachim, gest. 1704} vom 27.\;Dezember desselben Jahres bezeugt (\textit{LSB} III,\;8 N.\;321, S.\;819.3--7).\cite{02607}
Demgemäß ist zu vermuten, dass die Bemerkung N.\;\ref{56214} im Zeitraum zwischen 1693 und 1701 enstanden ist.
\pend%
\pstart%
Es ist jedoch nicht auszuschließen, dass Leibniz schon früher von Cassinis\indexN{\textls{Cassini}, Giovanni Domenico 1625\textendash1712} Beobachtungen und Schlussfolgerungen über die Zusammensetzung der Luft erfahren haben könnte, etwa während seines Aufenthaltes in Paris\indexO{Paris} 1672 bis 1676 oder ebenfalls zu einem späteren Zeitpunkt.
Demnach könnte N.\;\ref{56214} auch vor 1693 entstanden sein.
Die Genauigkeit, mit der Leibniz in N.\;\ref{56214} Cassinis\indexN{\textls{Cassini}, Giovanni Domenico 1625\textendash1712} Standpunkt wiedergibt, lässt aber eher mutmaßen, dass der Text erst nach Leibnizens Lektüre des besagten \textit{Recueil d’observations}\cite{02607} abgefasst wurde.
\pend%
\pstart%
Nach heutigem Wissensstand lässt sich ferner nicht feststellen, ob die Bemerkung N.\;\ref{56214} ursprünglich Teil eines eigenständigen Textes von Leibniz war oder vielmehr zu einem Auszug aus einem fremden Text gehörte.
\pend%
\end{ledgroup}%
\selectlanguage{latin}%
\frenchspacing%
\vspace{8mm}%
\newpage%
\normalsize%
\pstart%
\noindent%
\lbrack S.\;5\rbrack\
l'huyle%
\indexS{huyle}
qui est plu legere que l'eau%
\indexS{eau}
et aussi moins dense,
fait pourtant une moins grande refraction.%
\indexS{refraction}
Mons.
\edlabel{Gerland_1906_005b_Cassini_wmc-1}%
\edtext{%
Cassini dit,%
\indexN{\textls{Cassini}, Giovanni Domenico 1625\textendash1712}
que la hauteur de l'air%
\indexS{hauteur de l'air}%
\indexS{air}
trouvée par les refractions%
\indexS{refraction}
est moindre,
que celle
qui se trouve par les raisonnemens sur le Barometre.%
\indexS{barometre}%
}{%
\lemma{Cassini \lbrack...\rbrack\ Barometre}%
\Cfootnote{%
Vgl. \textsc{G.\;D. Cassini},
\textit{Observations astronomiques faites en divers endroits du royaume pendan l’année 1672} (\textsc{G.\;D. Cassini u. Andere}, \textit{Recueil d’observations faites en plusieurs voyages par ordre de sa majesté, pour perfectionner l’astronomie et la géographie}, Paris 1693, 4.\;Beitrag, S.\;20.)\cite{02605}%
}}%
\edlabel{Gerland_1906_005b_Cassini_wmc-2}
Ainsi
\edtext{%
il y a apparence,
qu'il y a deux matieres%
\indexS{matières qui composent l'air}
qui composent l'air:%
\indexS{air}%
}{%
\lemma{il y a \lbrack...\rbrack\ l'air}%
\Cfootnote{%
Vgl. a.a.O.\cite{02605}
}}
l'une
unie comme une huyle,%
\indexS{huyle}
qui est repercée dans certaines bornes
et qui fait la refraction%
\indexS{refraction}%
\lbrack;\rbrack\
l'autre est discontinue,
comme une matiere branchue %
\indexS{matière branchue}
qui passe au dessus de cette huyle%
\indexS{huyle}
et se etend par son ressort;%
\indexS{ressort de l'air}
c'est celle
qui fait la pesanteur de l'atmosphère%
\indexS{pesanteur de l'atmosphère}%
\indexS{atmosphère}
\edtext{%
\dots%
}{%
\lemma{\dots\,}%
\Cfootnote{%
Die Zeichensetzung in \textit{E} deutet darauf hin,
dass der Text unvollständig sein dürfte.%
}}
\pend%
\selectlanguage{latin}%
\frenchspacing%
\normalsize%
\newpage%
